\section{Teorema de Lucas}
\whenuse{Calcular $\binom nk\bmod p$ con $p$ primo peque\~no y $n,k$ demasiado
grandes para precomputar factoriales hasta $n$.}

Usar para calcular $\binom nk\bmod p$ cuando $p$ es primo peque\~no pero $n$ y
$k$ pueden ser enormes. Se descomponen ambos n\'umeros en base $p$.

\complexity{preproceso $O(p)$; consulta $O(\log_p n)$; memoria $O(p)$.}

\begin{lstlisting}[style=cpp]
struct LucasTheorem {
    ll prime;
    vector<ll> factorial, inverse_factorial;

    explicit LucasTheorem(ll modulus)
        : prime(modulus), factorial(prime), inverse_factorial(prime) {
        factorial[0] = 1;
        for (ll value = 1; value < prime; ++value)
            factorial[value] = factorial[value - 1] * value % prime;
        inverse_factorial[prime - 1] = mod_pow(factorial[prime - 1], prime - 2);
        for (ll value = prime - 1; value > 0; --value)
            inverse_factorial[value - 1] = inverse_factorial[value] * value % prime;
    }

    ll mod_pow(ll base, ll exponent) const {
        ll result = 1;
        while (exponent) {
            if (exponent & 1) result = result * base % prime;
            base = base * base % prime;
            exponent >>= 1;
        }
        return result;
    }

    ll small_combination(ll n, ll k) const {
        if (k < 0 || k > n) return 0;
        return factorial[n] * inverse_factorial[k] % prime
             * inverse_factorial[n - k] % prime;
    }

    ll combination(ll n, ll k) const {
        ll answer = 1;
        while (n || k) {
            answer = answer * small_combination(n % prime, k % prime) % prime;
            n /= prime; k /= prime;
        }
        return answer;
    }
};
\end{lstlisting}

\warning{Requiere $p$ primo y guardar $O(p)$ valores; no usar si el m\'odulo
es grande.}
